% analysis/pipeline_diagram.tex

\section{AIM Pipeline Overview}

This section answers a system-level question: which operator sequence
connects CRS context, intrinsic modelling, realization, and feedback
without breaking conceptual ownership across layers?

The figure is introduced as this operator-sequencing answer.

\begin{figure}[ht]
\centering
\begin{tikzpicture}[
  node distance=1.45cm,
  box/.style={
    draw,
    rectangle,
    rounded corners,
    minimum width=6.2cm,
    minimum height=0.95cm,
    align=center
  },
  arrow/.style={->, thick},
  feedback/.style={->, thick, dashed}
]

% Nodes
\node[box] (crs) {CRS Transformation \\ $\mathcal{C}$};

\node[box, below=of crs] (world) {World Coordinates};

\node[box, below=of world] (w2t) {World $\rightarrow$ Track \\ $\mathcal{W}$};

\node[box, below=of w2t] (model) {Curvature Model \\ $\kappa(s)$ \\ (parametric + hybrid)};

\node[box, below=of model] (real) {Realization \\ $\mathcal{R}$};

\node[box, below=of real] (t2w) {Track $\rightarrow$ World \\ $\mathcal{W}^{-1}$};

\node[box, below=of t2w] (output) {Realized Geometry};

% Main pipeline arrows
\draw[arrow] (crs) -- (world);
\draw[arrow] (world) -- (w2t);
\draw[arrow] (w2t) -- (model);
\draw[arrow] (model) -- (real);
\draw[arrow] (real) -- (t2w);
\draw[arrow] (t2w) -- (output);

% Feedback loops
\draw[feedback]
  (output.east) -- ++(1.2,0) |- node[pos=0.25, right]{Optimization / SQP} (model.east);

\draw[feedback]
  (output.west) -- ++(-1.2,0) |- node[pos=0.25, left]{Measurement / Data} (w2t.west);

\end{tikzpicture}

\caption{AIM processing pipeline as composition of operators linking
coordinate systems, alignment modelling, realization, and optimization.}
\end{figure}

Interpreted in chapter context, the diagram establishes a forward
processing spine with two explicit feedback channels: measurement data
stabilizes projection interpretation, and optimization feedback updates
the curvature model.

The engineering consequence is that pipeline correctness depends on
operator composition discipline, not on isolated local computations.
This prepares the following sections, where pipeline behaviour is
discussed through interaction patterns and failure modes.
